%% parts/p08/conditionals-deontic-free.tex --- Part 8
\chapter{Conditionals, Deontic, and Free Logics}
\label{sec:p08-conditionals-deontic-free}

\section{Free logic and non-denoting terms}

\begin{hbdefinition}{def:p08-negative-free-logic}{Negative free-logic semantics}
Fix a nonempty outer domain \(D^{\mathrm{out}}\), a nonempty existing domain
\(D^{\mathrm{in}}\subseteq D^{\mathrm{out}}\), and a total interpretation of
individual constants and function symbols into \(D^{\mathrm{out}}\).  In the
selected \emph{negative free-logic semantics}, each \(n\)-ary predicate has
extension contained in \((D^{\mathrm{in}})^n\),
\(E!(t)\) is true exactly when the value of \(t\) belongs to
\(D^{\mathrm{in}}\), and strict identity \(r=t\) is true exactly when both
term values exist and are the same object.  Quantifiers range over
\(D^{\mathrm{in}}\).  Hence every atom containing a denoting but nonexisting
outer-domain object is false; this is not called ``non-denotation.''

The adjective ``negative'' is necessary because positive free logic permits
some true predications of nonexisting objects, while a neutral partial
semantics can assign a truth-value gap when a term has no value.  It names one
two-valued model class, not all free logics.  Its strict atomic and equality
clauses distinguish it from the two exact variants in the table below and are
used by the guarded quantifier rules and description policy
\cite{Lambert1967,Bencivenga1986}.
\index{free logic}\index{negative free logic}\index{existence predicate}
\end{hbdefinition}

For a Russellian contextual treatment, an occurrence
\(B(\iota xA(x))\) is eliminated as
\[
 \exists x\bigl(A(x)\land
   \forall y(A(y)\to y=x)\land B(x)\bigr).
\]
In a primitive-description treatment, \(\iota xA(x)\) denotes the unique
member of \(D^{\mathrm{in}}\) satisfying \(A\), if there is one; otherwise the
term is undefined.  Under the primitive negative policy, every atom with that
undefined argument---including identity and \(E!\)---is false.  Under the
neutral policy it has value \(\mathsf u\), which then propagates by the
strong-Kleene operations.  These partial-description extensions do not alter
the total interpretation of the ordinary constants and function symbols fixed
above.  Russellian elimination and primitive descriptions are different formal
presentations, not two notations for one rule
\cite{Russell1905,Lehmann2002}.

\begin{table}[htbp]
\centering
\begin{tabularx}{\textwidth}{@{}l X X X@{}}
\toprule
Variant & Term and atomic semantics & Equality/quantifiers & Description policy \\
\midrule
negative & ordinary terms have total outer denotation; a failed primitive
description is undefined and makes its containing atom false; predicate
extensions lie in \((D^{\mathrm{in}})^n\) & strict
identity; quantifiers range over \(D^{\mathrm{in}}\) & Russell elimination or
the unique-existing primitive clause above \\
positive & total outer denotation; predicate extensions may contain tuples from
\((D^{\mathrm{out}})^n\) & outer equality is ordinary identity; quantifiers
range over \(D^{\mathrm{in}}\), and \(E!\) does not follow from arbitrary
predication & a primitive
description may denote an outer object only under a separately stated policy \\
neutral & terms may be partial; an atom with an undefined argument has value
\(\mathsf u\), and compounds use the strong-Kleene tables & identity is
\(\mathsf u\) with an undefined argument; only \(\mathsf t\) is designated;
quantifiers take the strong-Kleene infimum/supremum over
\(D^{\mathrm{in}}\) & a failed primitive description is undefined and hence
propagates \(\mathsf u\) at atomic occurrences \\
\bottomrule
\end{tabularx}
\caption{Free-logic variants with fixed denotation, atomic, identity,
quantifier, and description policies.}
\label{tab:p08-free-logic-variants}
\end{table}

\begin{hbdefinition}{def:p08-free-logic-consequence}{Variant-indexed free-logic consequence}
For \(v\in\{\mathrm{neg},\mathrm{pos},\mathrm{neu}\}\), write
\(\Gamma\models_v A\) when every model of the corresponding row of
Table~\ref{tab:p08-free-logic-variants} that designates every member of
\(\Gamma\) also designates \(A\).  In the negative and positive rows,
designation is ordinary truth; in the neutral row, only
\(\mathsf t\) is designated.  These are three typed semantic consequence
relations, not one family-wide proof calculus.  The subscripts are necessary
because the atomic and partial-term policies differ, and they are used in the
quantifier-rule table below.
\end{hbdefinition}

\begin{table}[htbp]
\centering\small
\begin{tabularx}{\textwidth}{@{}l X X X@{}}
\toprule
Variant & Universal instantiation & Existential generalization & Exact guard \\
\midrule
negative &
\(\forall x\,A(x),E!(t)\models_{\mathrm{neg}}A(t)\) &
\(A(t),E!(t)\models_{\mathrm{neg}}\exists x\,A(x)\) &
\(E!(t)\) must be true, so the outer-denoting term denotes an existing
inner-domain object \\
positive &
\(\forall x\,A(x),E!(t)\models_{\mathrm{pos}}A(t)\) &
\(A(t),E!(t)\models_{\mathrm{pos}}\exists x\,A(x)\) &
the same existence guard is required because positive atomic predication of an
outer object does not put that object in the quantifier domain \\
neutral &
\(\forall x\,A(x),E!(t)\models_{\mathrm{neu}}A(t)\) &
\(A(t),E!(t)\models_{\mathrm{neu}}\exists x\,A(x)\) &
both premises must have designated value \(\mathsf t\); an undefined or
nonexisting term gives \(E!(t)\) a nondesignated value and licenses no step \\
\bottomrule
\end{tabularx}
\caption{Quantifier-rule surfaces for the three fixed free-logic semantics.
The shared written guard has variant-specific denotation and consequence
standing; the table asserts no interderivability of proof calculi.}
\label{tab:p08-free-logic-quantifier-rules}
\end{table}

Consequently the classical theorem \(\exists x(x=t)\) is available only under
the corresponding existence/definedness guard
\cite{GratzlPavlovicNolt2025}.

\section{Deontic logic}

\begin{hbdefinition}{def:p08-sdl}{Standard deontic logic}
Standard deontic logic \(\mathsf{SDL}\) is normal modal logic \(\mathsf{KD}\)
with \(O A\) read as ``it ought to be that \(A\)'' and
\(\mathsf P_O A\defeq\neg O\neg A\). Relationally, \(O\) is box over a serial
ideal-alternative relation. The name is necessary to fix this monadic normal
system rather than deontic logic generally. It stands as a formal modal
calculus/semantics, differs through seriality and the D axiom
\(OA\to\mathsf P_OA\), and is used as the baseline for contrary-to-duty failures
\cite{VonWright1951,McNamara2006}.
\index{standard deontic logic}\index{obligation operator}\index{permission operator}
\end{hbdefinition}

Normality immediately yields the formal Ross inference
\[
 OA\quad\Longrightarrow\quad O(A\lor B),
\]
because \(A\to A\lor B\) is a theorem and necessitation plus distribution
preserves the implication.  Reading \(A\) as mailing a letter and \(B\) as
burning it makes the conclusion pragmatically weaker or misleading, but the
derivation is valid in SDL; the diagnosis concerns the adequacy of monadic
normal obligation for natural-language permission, not a failure of its rules.

The Good Samaritan pattern begins with \(O(H\land R)\), read as ``it ought to
be that a Samaritan helps someone who has been robbed.''  SDL derives \(OR\)
from the theorem \((H\land R)\to R\).  The unwanted reading that the robbery
itself ought to occur results from placing the descriptive circumstance inside
the scope of monadic \(O\).  A scope-sensitive or dyadic analysis treats
\(R\) as a condition instead.

For the Gentle Murder pattern, let \(K\) say that the agent kills and \(G\)
that the agent kills gently, with \(G\to K\).  From
\[
 O\neg K,\qquad K\to OG,\qquad K
\]
ordinary detachment gives \(OG\); normality and \(G\to K\) then give \(OK\),
contrary to the primary obligation \(O\neg K\).  The violated-condition norm
should therefore remain conditional, for example \(O(G\mid K)\), rather than
become an unqualified monadic obligation merely because \(K\) is factually
true \cite{Forrester1984}.
\index{Gentle Murder paradox}

For Chisholm's contrary-to-duty pattern, let \(A\) say that the agent goes to
help and \(B\) that the agent gives notice.  The set
\[
 OA,\qquad O(A\to B),\qquad \neg A\to O\neg B,\qquad \neg A
\]
detaches both \(OB\) (from the first two obligations) and \(O\neg B\) (from
the violated-condition rule).  Moving all conditionals inside \(O\), or all
outside it, changes the derivations but does not by itself distinguish primary
from violation-triggered obligations \cite{Chisholm1963,McNamara2006}.

\begin{hbdefinition}{def:p08-dyadic-obligation}{Best-world semantics for dyadic obligation}
A dyadic deontic model adds, at each world \(w\), a total preorder
\(\preceq_w\) on
\(R(w)=\{u\in W:wRu\}\), oriented so that
\(u\preceq_w v\) means that \(u\) is at least as ideal as \(v\).  Its strict
part is
\[
 u\prec_w v
 \quad\Longleftrightarrow\quad
 u\preceq_w v\ \text{and}\ v\not\preceq_w u.
\]
The preorder is \emph{smooth on definable candidate sets} when, for every
formula \(B\) and every \(u\in R(w)\cap[B]\), either no member of
\(R(w)\cap[B]\) is strictly more ideal than \(u\), or there is
\(v\in R(w)\cap[B]\) such that \(v\prec_w u\) and no
\(z\in R(w)\cap[B]\) satisfies \(z\prec_w v\).  Define, for every \(B\),
\[
 \operatorname{Best}_w(B)
 =\{u\in R(w)\cap[B]:
       \nexists v\in R(w)\cap[B]\ (v\prec_w u)\}.
\]
Thus \(\operatorname{Best}_w(B)=\varnothing\) when there is no accessible
\(B\)-world, while smoothness makes it nonempty whenever the candidate set is
nonempty.  Put
\[
 M,w\models O(A\mid B)
 \quad\Longleftrightarrow\quad
 \operatorname{Best}_w(B)\ne\varnothing
 \quad\text{and}\quad
 \operatorname{Best}_w(B)\subseteq[A].
\]
The existence conjunct makes an obligation with an impossible condition false
in this selected presentation; a vacuous alternative would be a different
dyadic semantics.
The name is necessary because \(B\) is a parameter of the ideality comparison,
not the antecedent of material implication.  It names a binary modality under
the displayed preference semantics, not dyadic deontic logic in general.  It
differs from \(O(B\to A)\) and is used to state
\(O(\neg B\mid\neg A)\) as a violation-triggered obligation without detaching
it in ideal \(A\)-worlds \cite{Hansson1969}.
\index{dyadic deontic logic}\index{conditional obligation}
\end{hbdefinition}

\begin{table}[htbp]
\centering\small
\begin{tabularx}{\textwidth}{@{}l X X@{}}
\toprule
Item & Formal surface & Exact lesson \\
\midrule
SDL & \(\mathsf{KD}\) with \(O\),
\(\mathsf P_OA=\neg O\neg A\) & serial monadic normal baseline \\
Ross & \(OA\vdash O(A\lor B)\) & valid SDL weakening; pragmatic permission
reading is not preserved \\
Good Samaritan & \(O(H\land R)\vdash OR\) & descriptive condition was placed
inside monadic obligation \\
Gentle Murder & \(O\neg K,K\to OG,K,G\to K\) & detachment and normality yield
both \(OK\) and \(O\neg K\); the violation norm needs conditional standing \\
Chisholm & \(OA,O(A\to B),\neg A\to O\neg B,\neg A\) & primary and
violation-triggered duties collapse under mixed monadic encodings \\
dyadic preference & \(O(A\mid B)\) at best accessible \(B\)-worlds & condition
selects the comparison class; no material-implication reduction \\
\bottomrule
\end{tabularx}
\caption{Deontic systems and paradox diagnoses.}
\label{tab:p08-deontic-systems}
\end{table}

Input/output logics and norm systems instead treat norms as pairs rather than
truth-valued formulas.  Modal truth, preference selection, and normative
closure have different standing and consequence relations.

\section{Counterfactual conditionals}

\begin{hbdefinition}{def:p08-stalnaker-conditional}{Selection-function counterfactual}
In a Stalnaker model, a selection function has type
\(f:W\times(\mathcal P(W)\setminus\{\varnothing\})\to W\) and satisfies
\(f(w,X)\in X\).  Put
\[
 M,w\models A>B
 \quad\Longleftrightarrow\quad
 [A]=\varnothing
 \quad\text{or}\quad
 M,f(w,[A])\models B.
\]
Thus impossible antecedents are vacuously true in this selected presentation;
for nonempty antecedents exactly one world is selected.  Centering is the
additional condition \(w\in X\Rightarrow f(w,X)=w\), not part of selection
alone.  The definition is needed because the
connective \(>\) is not material implication and because uniqueness is stronger
than merely having a set of closest worlds. It stands as a semantic clause,
differs from Lewis's sphere semantics through unique selection, and is used to
test conditional excluded middle \cite{Stalnaker1968}.
\index{counterfactual conditional}\index{selection-function semantics}
\end{hbdefinition}

\begin{hbdefinition}{def:p08-lewis-conditional}{Lewis sphere counterfactual}
A Lewis model assigns each world \(w\) a nonempty family
\(\mathcal S_w\subseteq\mathcal P(W)\) satisfying:
\begin{enumerate}
  \item \(\mathcal S_w\) is linearly ordered by inclusion;
  \item for every \(\mathcal T\subseteq\mathcal S_w\),
    \(\bigcup\mathcal T\in\mathcal S_w\); and
  \item if \(\varnothing\ne\mathcal T\subseteq\mathcal S_w\) and
    \(\bigcap\mathcal T\ne\varnothing\), then
    \(\bigcap\mathcal T\in\mathcal S_w\).
\end{enumerate}
Write \(U_w=\bigcup\mathcal S_w\).  The variably strict clause is
\[
 M,w\models A\Box\!\!\to B
 \quad\Longleftrightarrow\quad
 [A]\cap U_w=\varnothing
 \quad\text{or}\quad
 \exists S\in\mathcal S_w
 \bigl(S\cap[A]\ne\varnothing\land S\cap[A]\subseteq[B]\bigr).
\]
This clause does not require a least \(A\)-permitting sphere.  The system is
\emph{weakly centered} when, for every \(w\),
\[
 w\in U_w
 \quad\text{and}\quad
 \forall S\in\mathcal S_w\,
   (S\ne\varnothing\to w\in S),
\]
and \emph{centered} when
\(\{w\}\in\mathcal S_w\) for every \(w\).  Nesting makes centering imply weak
centering.  Lewis's \(\mathsf V\), \(\mathsf{VW}\), and \(\mathsf{VC}\)
use, respectively, arbitrary sphere systems of the displayed kind, weakly
centered systems, and centered systems.  The \emph{limit assumption} is the
further condition that, whenever \([A]\cap U_w\ne\varnothing\),
\[
 \exists S\in\mathcal S_w\,
 \bigl(S\cap[A]\ne\varnothing\ \land\
   \forall T\in\mathcal S_w\,
     (T\cap[A]\ne\varnothing\to S\subseteq T)\bigr).
\]
None of these centering clauses entails the limit assumption.  The sphere
definition is needed because there may be several equally close worlds or no
minimal sphere. It stands as a comparative-similarity semantics, differs from
unique selection, and is used for limit-assumption and centering distinctions
\cite{Lewis1973}.
\index{sphere semantics}\index{limit assumption}
\end{hbdefinition}

\begin{hbdefinition}{def:p08-ck-semantics}{CK selection and premise semantics}
A \emph{CK selection model} assigns each world \(w\) a function
\(F_w:\mathcal P(W)\to\mathcal P(W)\), with no requirement that
\(F_w(X)\subseteq X\), and sets
\[
 M,w\models A>B
 \quad\Longleftrightarrow\quad
 F_w([A])\subseteq[B].
\]
If \(F_w(\varnothing)=\varnothing\), every conditional with an impossible
antecedent is true at \(w\); basic CK does not impose that additional condition.
Equivalently, a \emph{premise-semantics presentation} assigns a relation
\(\mathcal N_w\subseteq\mathcal P(W)\times\mathcal P(W)\) and reads
\(A>B\) as \(([A],[B])\in\mathcal N_w\), where each fixed-antecedent section
is a principal filter: for some core \(K_w(X)\subseteq W\),
\((X,Y)\in\mathcal N_w\) exactly when \(K_w(X)\subseteq Y\).  Such a section
contains \(W\) and is closed under intersection and supersets.  The two
presentations correspond by
\[
 (X,Y)\in\mathcal N_w
 \quad\Longleftrightarrow\quad F_w(X)\subseteq Y.
\]
The names are necessary because assigning a consequent-test set to an
antecedent proposition and relating premise/conclusion propositions are
different semantic data.  They
name model presentations for the same basic logic, not object-language
connectives.  Their nonredundancy is visible in the displayed conversion, and
the CK completeness theorem uses it \cite{Chellas1975}.  The absent inclusion
condition is load-bearing: adding \(F_w(X)\subseteq X\) validates the identity
schema \(A>A\), which belongs to stronger conditional systems rather than
basic \(\mathsf{CK}\).
\index{conditional logic!CK}\index{premise semantics!conditional logic}
\end{hbdefinition}

\begin{hbtheorem}{thm:p08-ck-completeness}{cited=Chellas1975}{Soundness and completeness of conditional logic CK}
The theoremhood calculus on
\hyperref[sheet:hilbert-conditional-ck]{\texttt{sheet:hilbert-conditional-ck}}
is sound and complete for the CK selection models, equivalently the
fixed-antecedent premise semantics, of
Definition~\ref{def:p08-ck-semantics}.
\end{hbtheorem}

\begin{hbtheorem}{thm:p08-vc-completeness}{cited=Lewis1971}{Soundness and completeness of Lewis VC, or C1}
Lewis's centered variably strict system \(\mathsf{VC}\), called
\(\mathsf{C1}\) in the cited alternative nomenclature, is presented on
\hyperref[sheet:hilbert-counterfactual-vc]{\texttt{sheet:hilbert-counterfactual-vc}}
and is sound and complete for the centered sphere semantics of
Definition~\ref{def:p08-lewis-conditional}.  The source's limit-free clause is
the one printed there; no least-sphere assumption is smuggled into the theorem
\cite{Wen2019}.
\end{hbtheorem}

\begin{systementry}{entry:p08-conditional-ck}{Chellas basic conditional logic CK}
  \entrylanguage{Classical propositional logic with primitive binary
    conditional \(A>B\).}
  \entrysemantics{Set-valued selection functions, equivalently principal
    premise semantics, Definition~\ref{def:p08-ck-semantics}.}
  \entrycalculus{sheet:hilbert-conditional-ck}
  \entryadequacy{cited=Chellas1975}{Sound and complete for the exact CK model
    class, Theorem~\ref{thm:p08-ck-completeness}.}
  \entrydecidability{cited=Nute1980}{The propositional theoremhood problem is
    decidable for the displayed finitary language.}
  \entrycharacteristic{RCEA and RCK govern antecedent extensionality and normal
    reasoning inside the consequent; centering and CEM are absent.}
  \entrytranslations{Principal premise sections correspond exactly to their
    selection cores by Definition~\ref{def:p08-ck-semantics}.}
  \entrysources{Chellas1975,Nute1980}
\end{systementry}

\begin{systementry}{entry:p08-counterfactual-vc}{Lewis counterfactual logic VC, or C1}
  \entryaliases{\(\mathsf{C1}\) in the cited alternative nomenclature.}
  \entrylanguage{The CK conditional language.}
  \entrysemantics{Centered nested sphere systems under the exact limit-free
    variably strict clause of Definition~\ref{def:p08-lewis-conditional}.}
  \entrycalculus{sheet:hilbert-counterfactual-vc}
  \entryadequacy{cited=Lewis1971}{Sound and complete for centered sphere
    semantics, Theorem~\ref{thm:p08-vc-completeness}.}
  \entrydecidability{cited=Lewis1971}{The propositional theoremhood problem is
    decidable.}
  \entrycharacteristic{ID, CA, CSO, CV, conditional MP, and centering;
    unique selection and CEM are not assumed.}
  \entrytranslations{The flat fragment connects to preferential reasoning only
    under its stated restriction; no unrestricted identification is made.}
  \entrysources{Lewis1971,Lewis1973,Wen2019}
\end{systementry}

\begin{table}[htbp]
\centering\scriptsize
\begin{tabularx}{\textwidth}{@{}l X X X@{}}
\toprule
System/presentation & Exact model surface & Characteristic principles & Boundary \\
\midrule
\(\mathsf{CK}\) & set-valued selection or equivalent premise semantics &
RCEA and RCK over CPL & no centering, uniqueness, ID, or CEM assumed;
antecedent strengthening and conditional transitivity fail in general \\
Lewis \(\mathsf V\) & nested sphere systems & ID, CA, CSO, CV & no centering \\
Lewis \(\mathsf{VW}\) & weakly centered spheres & \(\mathsf V\) plus
conditional MP & actual world is among closest antecedent worlds \\
Lewis \(\mathsf{VC}=\mathsf{C1}\) & centered spheres &
\(\mathsf{VW}\) plus CS & multiple equally close antecedent worlds permitted;
antecedent strengthening and conditional transitivity fail in general \\
Stalnaker \(\mathsf{C2}\) & unique selection with centering & CEM
\((A>B)\lor(A>\neg B)\) in addition to the centered core & uniqueness and the
impossible-antecedent policy are load-bearing; antecedent strengthening and
conditional transitivity still fail in general \\
\bottomrule
\end{tabularx}
\caption{Conditional-logic family map: CK, Lewis's variably strict neighbors,
and Stalnaker's unique-selection extension.}
\label{tab:p08-conditional-logics}
\end{table}

The two commonly invalid schemas are
\[
 \mathrm{AS}\colon (A>B)\to((A\land C)>B),
 \qquad
 \mathrm{Trans}\colon ((A>B)\land(B>C))\to(A>C).
\]
Both failures already occur in centered Stalnaker models.  For AS, take
\([A]=\{u,v\}\), \([B]=\{u\}\), \([C]=\{v\}\), let the actual world satisfy
none of \(A,B,C\), and set \(f(w,[A])=u\) and
\(f(w,[A\land C])=v\).  For Trans, take
\([A]=\{u\}\), \([B]=\{u,v\}\), \([C]=\{v\}\) and set
\(f(w,[A])=u\), \(f(w,[B])=v\).  Complete the selection function
arbitrarily subject to selection and centering.  The first model makes AS
false; the second makes both premises of Trans true and its conclusion false.
Thus counterfactuals also generally fail contraposition and monotonicity in the
antecedent. Modus ponens requires a centering or actuality condition.
Conditional excluded middle
\((A>B)\vee(A>\neg B)\) is validated by unique selection but may fail when
several closest \(A\)-worlds disagree. Proof systems must therefore name the
semantic conditions they capture \cite{Nute1980,Bennett2003}.

\section{Indicative conditionals}

Material implication validates both
\[
 \neg A\vdash A\to B,
 \qquad
 B\vdash A\to B.
\]
These are the paradoxes of material implication in the present comparison:
the falsity of an antecedent or independent truth of a consequent suffices for
the truth-functional conditional, even when it supplies no inferential or
assertability connection between \(A\) and \(B\).

\begin{hbdefinition}{def:p08-adams-conditional}{Probabilistic assertability of an indicative}
For events with \(P(A)>0\), the Adams-style assertability of ``if \(A\), then
\(B\)'' is governed by \(P(B\mid A)\). The terminology is necessary because a
conditional probability is not a truth value assigned by the material
conditional. It stands as an assertability/uncertainty measure, differs from a
truth-functional connective, and is used to explain why high
\(P(\neg A\vee B)\) caused merely by low \(P(A)\) need not support the
indicative \cite{Adams1975}.
\index{indicative conditional}\index{conditional probability}
\end{hbdefinition}

\begin{hbtheorem}{thm:p08-lewis-triviality}{cited=Lewis1976}{A probability-of-conditionals triviality boundary}
Let \(\mathcal B\) be the finite Boolean algebra generated by logically
independent events \(A,B\), so its four \(A,B\)-atoms are nonempty.  There is no
fixed event \(C\in\mathcal B\) such that
\[
 P(C)=P(B\mid A)
\]
for every strictly positive probability \(P\) on \(\mathcal B\).
Finiteness (or an explicit measure-rich extension hypothesis), the fixed-event
assumption, and membership in the same algebra are essential.  This is the
finite four-atom consequence of Lewis's source-carried triviality result; no
claim about a proof internal to this handbook is made.
\end{hbtheorem}

The theorem constrains one proposition-plus-probability analysis; it does not
show that indicatives lack semantics.  In a \emph{suppositional account}, an
information state is first restricted or temporarily updated by \(A\), and
\(B\) is then assessed in that updated state.  This policy names an
assessment procedure rather than an event \(C\) in the original algebra, so it
does not assert the identity excluded by the theorem.  The name is necessary
to distinguish temporary antecedent supposition from truth-functional
evaluation and fixed-event representation.  Dynamic,
relevance-sensitive, and inferential versions require their own update and
acceptance rules \cite{Jackson1987,Edgington1995,Bennett2003}.

\begin{hbwarning}{def:p08-conditional-boundary}{One glyph does not determine one conditional}
Material implication, strict implication, Stalnaker selection,
Lewis spheres, dyadic obligation, and conditional probability have different
truth or assertability clauses. A theorem transfers between them only through
an explicit translation with proved preservation; ordinary-language ``if''
does not provide that proof.
\end{hbwarning}
